\documentclass[12pt,a4paper]{article}

% -------------------------------------------------
% Sprache / Kodierung
% -------------------------------------------------
\usepackage[T1]{fontenc}
\usepackage[utf8]{inputenc}
\usepackage[ngerman]{babel}
\usepackage{lmodern}
\usepackage{microtype}

% -------------------------------------------------
% Mathematik
% -------------------------------------------------
\usepackage{amsmath,amssymb,amsthm,mathtools}

% -------------------------------------------------
% Layout / Grafik / Farben
% -------------------------------------------------
\usepackage[a4paper,margin=2.3cm]{geometry}
\usepackage{booktabs}
\usepackage{xcolor}
\usepackage[hidelinks,unicode]{hyperref}
\pdfstringdefDisableCommands{%
  \def\ü{ue}\def\ö{oe}\def\ä{ae}\def\Ü{Ue}\def\Ö{Oe}\def\Ä{Ae}\def\ss{ss}%
}
\usepackage[most]{tcolorbox}
\usepackage{enumitem}
\usepackage{array}

% -------------------------------------------------
% Farben
% -------------------------------------------------
\definecolor{lückenrot}{RGB}{180,40,40}
\definecolor{bewiesengrün}{RGB}{30,100,50}
\definecolor{warnorange}{RGB}{200,100,10}
\definecolor{hintergrundblau}{RGB}{235,242,250}
\definecolor{hintergrundgelb}{RGB}{255,252,225}
\definecolor{adischviolett}{RGB}{90,20,140}
\definecolor{fehlerrot}{RGB}{200,0,0}
\definecolor{konjviolett}{RGB}{80,0,130}
\definecolor{e8grau}{RGB}{55,55,90}
\definecolor{kingmanblau}{RGB}{20,60,130}
\definecolor{zirkelrot}{RGB}{160,0,0}
\definecolor{epiloggrau}{RGB}{45,45,65}
\definecolor{frontlinienrot}{RGB}{150,20,20}
\definecolor{neugrün}{RGB}{15,80,35}

% -------------------------------------------------
% tcolorbox-Stile
% -------------------------------------------------
\tcbset{
  lücke/.style={colback=red!8, colframe=lückenrot, fonttitle=\bfseries,
    title={Offene Lücke}},
  ergebnis/.style={colback=green!7, colframe=bewiesengrün,
    fonttitle=\bfseries, title={Bewiesenes Ergebnis}},
  warnung/.style={colback=orange!10, colframe=warnorange,
    fonttitle=\bfseries, title={Achtung / Heuristik}},
  kernidee/.style={colback=hintergrundblau, colframe=blue!60!black,
    fonttitle=\bfseries, title={Kernidee}},
  konj/.style={colback=violet!5, colframe=konjviolett, fonttitle=\bfseries,
    title={Konjektur (unbewiesen)}},
  epilog/.style={colback=gray!6, colframe=epiloggrau,
    fonttitle=\bfseries\large, title={Epilog}},
  frontlinie/.style={colback=red!5, colframe=frontlinienrot,
    fonttitle=\bfseries, title={Die Frontlinie}},
  saeulen/.style={colback=blue!4, colframe=kingmanblau,
    fonttitle=\bfseries, title={Die vier Säulen}},
  neufeld/.style={colback=green!5, colframe=neugrün,
    fonttitle=\bfseries, title={Neue Forschungsrichtung}},
  adisch/.style={colback=violet!6, colframe=adischviolett,
    fonttitle=\bfseries, title={2-adische Perspektive}},
}

% -------------------------------------------------
% Theorem-Umgebungen
% -------------------------------------------------
\newtheorem{definition}{Definition}[section]
\newtheorem{proposition}[definition]{Proposition}
\newtheorem{theorem}[definition]{Theorem}
\newtheorem{lemma}[definition]{Lemma}
\newtheorem{korollar}[definition]{Korollar}
\newtheorem{bemerkung}[definition]{Bemerkung}
\newtheorem{hypothese}[definition]{Hypothese}
\newtheorem{satz}[definition]{Satz}

\newenvironment{beweis}{\begin{proof}[Beweis]}{\end{proof}}

% -------------------------------------------------
% Makros
% -------------------------------------------------
\newcommand{\vii}{\nu_2}
\newcommand{\N}{\mathbb{N}}
\newcommand{\Z}{\mathbb{Z}}
\newcommand{\R}{\mathbb{R}}
\newcommand{\Q}{\mathbb{Q}}
\newcommand{\H}{\mathbb{H}}
\newcommand{\Ztwo}{\mathbb{Z}_2}
\newcommand{\abs}[1]{\lvert#1\rvert}
\newcommand{\atwo}[1]{\lvert#1\rvert_2}
\newcommand{\Lam}{\Lambda}
\newcommand{\col}{\mathrm{Col}}
\newcommand{\KC}{K}
\newcommand{\TV}{\mathrm{TV}}

% -------------------------------------------------
% Dokument
% -------------------------------------------------
\title{%
  Die Collatz-Vermutung: Eine Bilanz\\[0.3em]
  \large Schlussartikel --- Grenzbefestigung, Beweislandschaft,
  offene Fragen\\[0.5em]
  \normalsize Zusammenfassung der Arbeiten 2025--2026%
}
\author{Thomas Hoffbauer}
\date{16.\ Juni 2026}

\begin{document}
\maketitle

\begin{abstract}
Dieser Schlussartikel zieht Bilanz über eine umfangreiche Untersuchung
der Collatz-Vermutung, die 16 LaTeX-Dokumente und 12 Pull-Requests
(davon 11 zur Grenzbefestigung) umfasst. Die vier zentralen Ergebnisse
sind: (1) ein rigoroser DC-Beweis für spezielle Klassen mit Lean-Formalisierung;
(2) der Kingman-Lyapunov-Exponent $\Lambda \approx -0{,}830$;
(3) die RPF-Mischungsrate des mod-12-Transfer-Operators; und
(4) die exakte Lokalisierung der verbleibenden Lücke als
Uniformitätsproblem ($\N$ als Nullmenge in $\Ztwo$).

\textbf{Die Collatz-Vermutung bleibt unbewiesen.} Dieser Artikel
dokumentiert, was bekannt ist, was bewiesen wurde, und -- mit dem
Epilog von Juni 2026 -- die neue Forschungsagenda der
\emph{arithmetischen Dynamik}.
\end{abstract}

\tableofcontents

\bigskip
\begin{tcolorbox}[warnung, title={Zum Status dieses Dokuments}]
Dieser Artikel ist ein Überblick und Abschlussbericht. Er setzt die
Kenntnis der Einzeldokumente
\textit{collatz\_uniformitaet.tex}, \textit{collatz\_lte\_dichte.tex},
\textit{collatz\_kingman.tex} und \textit{collatz\_arithmetische\_dynamik.tex}
voraus. Alle nicht-standardmäßigen Behauptungen sind in diesen Quellen
belegt oder dort als offen markiert.
\end{tcolorbox}

% ====================================================
\section{Die Kernfrage und der Stand der Forschung}
\label{sec:kernfrage}
% ====================================================

\subsection{Die Collatz-Vermutung}

\begin{definition}[Collatz-Abbildung]
Für $n \in \N$ sei:
\[
  T(n) :=
  \begin{cases}
    n/2 & \text{falls } n \equiv 0 \pmod{2}, \\
    (3n+1)/2 & \text{falls } n \equiv 1 \pmod{2}.
  \end{cases}
\]
Die \emph{Collatz-Vermutung} besagt: Für jedes $n \in \N$ gibt es
$k \in \N$ mit $T^k(n) = 1$.
\end{definition}

\begin{tcolorbox}[ergebnis, title={Taos Hauptergebnis (Forum Math.\ Pi, 2022)}]
Für jede Funktion $f: \N \to \R$ mit $f(n) \to \infty$ gilt:
\[
  \mathrm{Dens}_{\log}\bigl(\{n \in \N : \min_{k\geq 0}\col^k(n)
  \leq f(n)\}\bigr) = 1.
\]
Fast alle $n \in \N$ (logarithmisch) haben Collatz-Bahnen, die
beliebig klein werden.
\end{tcolorbox}

\subsection{Die Maß-Null-Barriere}

\begin{tcolorbox}[lücke, title={Das fundamentale Hindernis}]
$\N \subset \Ztwo^*$ ist eine abzählbare Menge, also eine $\mu_H$-Nullmenge
im 2-adischen Maß. Taos Ergebnis ($\mu_H(E) = 0$ für das Ausnahmeset $E$)
ist daher \emph{kein} Widerspruch zu $E \cap \N \neq \emptyset$.

\medskip
Die Collatz-Vermutung fordert ein \emph{individuelles} Resultat für jedes
$n \in \N$, während Maß\-argu\-mente nur \emph{generische} Aussagen liefern.
Das ist die Uniformitätslücke.
\end{tcolorbox}

% ====================================================
\section{Vier Säulen der Grenzbefestigung}
\label{sec:saeulen}
% ====================================================

\subsection{Säule 1: Der DC-Beweis}

Der \emph{Direct Contraction}-Beweis zeigt, dass bestimmte Klassen
von $n \in \N$ direkt konvergieren:

\begin{tcolorbox}[ergebnis]
Für ungerades $n \equiv 1 \pmod{4}$ gilt $T(n) = (3n+1)/2 < n$, da
$3n+1 \equiv 2 \pmod{4}$ und $T(n) = (3n+1)/2 < 2n$. Für $n \geq 3$:
$T(n) < n$. Lean-bewiesen.
\end{tcolorbox}

\subsection{Säule 2: Lyapunov-Exponent $\Lambda \approx -0{,}830$}

\begin{tcolorbox}[saeulen, title={Kingman-Lyapunov-Exponent}]
Aus dem subadditiven Kingman-Ergodensatz, angewandt auf die Collatz-Dynamik,
folgt (unter der Voraussetzung, dass die mod-12-Markov-Kette ergodisch ist):
\[
  \Lambda = \lim_{K\to\infty} \frac{1}{K}\sum_{k=0}^{K-1}
  \log_2\frac{T^{k+1}(n)}{T^k(n)} \approx -0{,}830.
\]
Das quantifiziert, was in Taos Ergebnis implizit liegt: Die durchschnittliche
Kontraktion pro Schritt ist $2^{\Lambda} \approx 0{,}562$.
\end{tcolorbox}

\subsection{Säule 3: RPF-Mischung des mod-12-Transfer-Operators}

\begin{tcolorbox}[saeulen, title={Exponentielle Mischungsrate}]
Die mod-12-Markov-Kette auf $\{E, A, B, C\}$ ist primitiv mit zweitem
Eigenwert $|\lambda_2| \approx 0{,}35$. Für alle $n \in \N$ und alle
$k \geq 0$:
\[
  \bigl\|M^k e_{B(n)} - \pi\bigr\|_{\TV} \leq C\,|\lambda_2|^k,
\]
wobei $C > 0$ nur von der mod-12-Startklasse abhängt. Die Mischungszeit
bis $\varepsilon = 0{,}01$ beträgt $\approx 5$ Blockschritte ---
\emph{unabhängig von $n$}.
\end{tcolorbox}

\subsection{Säule 4: Lean-Fundament}

Das Lean-Fundament umfasst:
\begin{itemize}[nosep]
  \item \texttt{padicValNat} für LTE-Berechnungen ($\vii_2(3^k - 1)$).
  \item Stirling-Approximation für asymptotische Abschätzungen.
  \item von-Staudt-Clausen für Bernoulli-Nenner.
  \item Vierquadrate-Satz als Referenz für Norm-Identitäten.
  \item DC-Beweis in Lean vollständig formalisiert.
\end{itemize}

% ====================================================
\section{Was nicht bewiesen werden konnte}
\label{sec:offen}
% ====================================================

\begin{tcolorbox}[warnung, title={Ehrliche Bilanz der Grenzen}]
\begin{enumerate}[label=(\arabic*), nosep, itemsep=3pt]
  \item \textbf{Schritt C (Quaternionen-Diophantik)} enthält eine
    versteckte Zirkularität: Ausnahmebahnen in $\N$ bleiben ganzzahlig,
    es entsteht kein Widerspruch.
  \item \textbf{Argument D (Binärinformation)} ist korrekt nur bedingt
    auf Nicht-Divergenz --- zirkulär.
  \item \textbf{Das Uniformitätsproblem} ist ungelöst: Die Brücke von
    ``$\mu_H(E) = 0$'' zu ``$E \cap \N = \emptyset$'' fehlt.
  \item \textbf{Spektrallücke} der mod-12-Kette ist numerisch belegt,
    aber nicht formal bewiesen.
  \item \textbf{Zweiblockkompensation} ist uniform bekannt, aber noch
    nicht für alle $(k,r)$ rigoros.
\end{enumerate}
\end{tcolorbox}

% ====================================================
% EPILOG: Die Frontlinie (Juni 2026)
% ====================================================

\section*{Epilog: Die Frontlinie (Juni 2026)}
\label{sec:epilog}
\addcontentsline{toc}{section}{Epilog: Die Frontlinie (Juni 2026)}

\begin{tcolorbox}[epilog]
Dieser Epilog wurde im Juni 2026 als Abschluss einer intensiven
Auseinandersetzung mit der Collatz-Vermutung geschrieben, die
\textbf{16 LaTeX-Dokumente} und \textbf{12 Pull-Requests} (PR~\#1
bis PR~\#12) umfasste. Der Zweck: die exakte verbleibende Frontlinie
zu lokalisieren --- nicht den Beweis zu vollenden, sondern den Abgrund
präzise zu kartieren.
\end{tcolorbox}

\subsection*{Was in dieser Arbeit geleistet wurde}

\begin{tcolorbox}[saeulen, title={16 Dokumente, 12 PRs: Die Grenzbefestigung}]
\begin{center}
\renewcommand{\arraystretch}{1.4}
\begin{tabular}{>{\bfseries}l >{\raggedright\arraybackslash}p{9cm}}
\toprule
PR \#1--3 & Grundlagen: DC-Beweis, LTE-Formel, EABC-Klassifikation.
            Lean-Formalisierungen. \\[2pt]
PR \#4--6 & Lyapunov-Analyse: Kingman-Exponent $\Lambda \approx -0{,}830$,
            Stirling, Bernoulli-Nenner. \\[2pt]
PR \#7--8 & LTE-Dichte: Extremale C-Ketten $n = 2^{k+1}3^r - 1$,
            Zweiblockkompensation. \\[2pt]
PR \#9    & Mathlib-Werkzeugkasten: \texttt{padicValNat}, Vierquadrate,
            von-Staudt-Clausen. \\[2pt]
PR \#10   & Kingman-Formalisierung: $\Lambda \approx -0{,}830$ als
            quantifizierter Lyapunov-Exponent. \\[2pt]
PR \#11   & Uniformitätsanalyse: RPF solide (Schritt B), Zirkularität
            von Schritt C und D dokumentiert. \\[2pt]
PR \#12   & \textbf{Arithmetische Dynamik (dieses Dokument):}
            Kolmogorov-Komplexität, Uniformitäts-Vermutung,
            Analogon zu Roths Satz. \\[2pt]
\bottomrule
\end{tabular}
\end{center}
\end{tcolorbox}

\subsection*{Die vier Säulen im Rückblick}

\begin{tcolorbox}[saeulen, title={Vier Säulen --- was steht, was wackelt}]
\begin{enumerate}[label=\textbf{Säule \arabic*:}, leftmargin=*, align=left,
    itemsep=4pt]
  \item \textbf{DC-Beweis (steht fest).} Für $n \equiv 1 \pmod{4}$: direkte
    Kontraktion $T(n) < n$. Lean-bewiesen. Aber nur ein Bruchteil aller $n$.

  \item \textbf{Lyapunov $\Lambda \approx -0{,}830$ (steht, bedingt).}
    Im Mittel schrumpft die Bahn um Faktor $2^\Lambda \approx 0{,}56$ pro
    Schritt. Bedingt auf RPF-Ergodizität --- die Brücke zu konkreten $n$
    fehlt.

  \item \textbf{RPF-Mischung (steht als Schranke).} Die mod-12-Kette mischt
    in $\leq 5$ Blockschritten, unabhängig von $n$. Formal für den Markov-Anteil
    korrekt. Der Übergang vom Markov-Modell zur echten Collatz-Dynamik ist
    das Uniformitätsproblem.

  \item \textbf{Lean-Fundament (steht als Werkzeugkasten).} Alle relevanten
    Mathlib-Primitiven sind identifiziert und einsatzbereit. Ein vollständiger
    Lean-Beweis der Collatz-Vermutung liegt nicht vor.
\end{enumerate}
\end{tcolorbox}

\subsection*{Die verbleibende Lücke: $\N$ als Nullmenge in $\Ztwo$}

\begin{tcolorbox}[frontlinie, title={Die Frontlinie: Eine Gleichung, eine Brücke}]
Die gesamte Arbeit läuft auf eine einzige, präzise formulierbare Lücke hinaus:

\medskip
\textbf{Bekannt:}
\[
  \mu_H\bigl(\{x \in \Ztwo^* : \text{Bahn von } x \text{ konvergiert nicht}\}\bigr) = 0
  \quad\text{(Tao 2022).}
\]

\textbf{Gesucht:}
\[
  \{n \in \N : \text{Bahn von } n \text{ konvergiert nicht}\} = \emptyset.
\]

\medskip
Das Problem: $\N$ ist eine $\mu_H$-Nullmenge in $\Ztwo$. Ein Maß-Null-Ergebnis
über $\Ztwo$ sagt \emph{nichts} über die $\mu_H$-Nullmenge $\N$. Das ist
nicht eine technische Kleinigkeit --- es ist der Abgrund.

\medskip
Alle Brückenversuche dieser Arbeit (Schritt C: Quaternionen-Zirkularität;
Argument D: Binärinformations-Zirkularität) sind an exakt diesem Abgrund
gescheitert. Die RPF-Mischung (Schritt B) ist der einzige Beitrag, der
nicht zirkulär ist --- aber er liefert auch keine Brücke, nur eine
quantitative Schranke.
\end{tcolorbox}

\subsection*{Die neue Forschungsfrage: Arithmetische Dynamik}

\begin{tcolorbox}[neufeld, title={Das neue Forschungsfeld (PR \#12)}]
Die Arbeit des zwölften Pull-Requests öffnet ein neues Forschungsfeld:
die \emph{arithmetische Dynamik} für Collatz.

\medskip
\textbf{Die zentrale neue Idee:}
Statt die Lücke direkt zu schließen, frage man nach der \emph{quantitativen
Trennung} zwischen $\N$ und dem Ausnahmeset $E$:
\[
  \mathrm{dist}_{2\text{-adisch}}(n, E) \geq c \cdot 2^{-\lfloor\log_2 n\rfloor}
  \quad \text{für alle } n \in \N \text{, mit universeller Konstante } c > 0.
\]

\medskip
\textbf{Das Analogon zu Roths Satz:} In der diophantischen Approximation
verhindert Roths Satz (1955), dass algebraische Irrationale durch Rationale
``zu gut'' approximiert werden. Das Collatz-Analogon würde verhindern, dass
natürliche Zahlen den hochkomplexen Fixpunkten von $T$ in $\Ztwo$ zu
nahe kommen.

\medskip
\textbf{Warum die Analogie vielversprechend ist:}
\begin{enumerate}[label=(\arabic*), nosep, itemsep=2pt]
  \item $\N$ hat Kolmogorov-Komplexität $O(\log n)$ --- endlich.
  \item Elemente von $E$ (falls $E \neq \emptyset$) haben generisch
    unendliche Kolmogorov-Komplexität.
  \item Der 2-adische Abstand zwischen endlich-komplexen und
    unendlich-komplexen Elementen ist messbar.
  \item Die LTE-extremalen Punkte $n = 2^{k+1}3^r - 1$ haben Abstand
    $\approx 2^{-(k+1)}$ zu benachbarten Zweierpotenzen --- konsistent
    mit der Schranke für $c = 1$.
\end{enumerate}

\medskip
\textbf{Was beweisbar sein müsste:}
Eine Version von Roths Satz für den Collatz-Operator auf $\Ztwo$ ---
eine Abstandsschranke aus der algebraischen Natur der Dynamik.
Dieses Ziel liegt außerhalb bekannter Methoden, ist aber präzise
genug formuliert, um als echter Forschungsgegenstand zu dienen.
\end{tcolorbox}

\subsection*{Bilanz: Was bleibt}

\begin{tcolorbox}[adisch, title={Abschlussbilanz der Grenzbefestigung}]
\textbf{Was diese Arbeit erreicht hat:}
\begin{itemize}[nosep, itemsep=2pt]
  \item Die Beweislandschaft wurde vollständig kartiert.
  \item Alle bekannten Argumente wurden auf Zirkularität geprüft.
  \item Die Lücke wurde auf eine einzige, präzise formulierbare Frage
    reduziert.
  \item Ein neues Forschungsfeld --- arithmetische Dynamik mit
    diophantischen Methoden --- wurde eröffnet.
  \item Der Lean-Werkzeugkasten ist bereit für eine zukünftige
    Formalisierung.
\end{itemize}

\medskip
\textbf{Was bleibt:}
Die Collatz-Vermutung. Der Abgrund zwischen Maß-Null-Argumenten
und individueller Konvergenz. Die Brücke zwischen $\mu_H(E) = 0$
und $E \cap \N = \emptyset$ ist nicht gebaut --- aber jetzt weiß man
genau, wie tief die Schlucht ist und welche Brücke fehlt:
eine Theorie der arithmetischen Dynamik, inspiriert von Roth.
\end{tcolorbox}

\subsection*{Referenzen}

\begin{itemize}
  \item T.\ Tao: \textit{Almost all orbits of the Collatz map attain
    almost bounded values}. Forum Math.\ Pi \textbf{10} (2022), e12.
  \item J.\ F.\ C.\ Kingman: \textit{The ergodic theory of subadditive
    stochastic processes}. J.\ Roy.\ Stat.\ Soc.\ B \textbf{30} (1968), 499--510.
  \item K.\ F.\ Roth: \textit{Rational approximations to algebraic numbers}.
    Mathematika \textbf{2} (1955), 1--20.
  \item D.\ Ruelle: \textit{Thermodynamic Formalism}. Addison-Wesley, 1978.
  \item Zugehörige Dokumente (PR \#1--\#12):
    \textit{collatz\_uniformitaet.tex},
    \textit{collatz\_lte\_dichte.tex},
    \textit{collatz\_arithmetische\_dynamik.tex}.
\end{itemize}

\end{document}
